\documentclass[documento]{gpt}
\newcommand{\subsubsubsection}[1]{\paragraph{#1}}
\setcounter{secnumdepth}{4} 
\setcounter{tocdepth}{4}
\usepackage{titlesec}
\titleformat{\paragraph}
  {\normalfont\normalsize\bfseries}
  {\theparagraph}{1em}{}
\titlespacing*{\paragraph}
  {0pt}{3.25ex plus 1ex minus .2ex}{1.5ex plus .2ex}


% ---------------------------------------------------------------------
% Dati specifici del documento
% ---------------------------------------------------------------------
\titolo{Piano di progetto}
\uso{esterno}
\destinatari{Prof.~Tullio Vardanega\newline Prof.~Riccardo Cardin}

% ---------------------------------------------------------------------
% Registro delle modifiche
% ---------------------------------------------------------------------
% Sintassi: \modifica{versione}{AAAA-MM-GG}{autore}{verificatore}{descrizione}
% La riga più recente va messa in cima.
%  Quando il Responsabile approva, aggiunge lui IN CIMA questa riga:
%  \modifica{1.0.0}{2026-10-23}{corro}{--}{Approvazione}
\modifica{0.2.0}{2026-11-04}{spagnoli, guglielmo}{corro}{Stesura della sezione Pianificazione e dello sprint 1}
\modifica{0.1.0}{2026-11-04}{spagnoli, guglielmo}{corro}{Prima stesura del documento}



% Scheda di un rischio:
% \rischio{codice}{titolo}{probabilità}{impatto}{livello}{descrizione}{rilevamento}{mitigazione}{contingenza}
\newcommand{\rischio}[9]{%
  \subsubsection{#1 -- #2}\label{rischio:#1}
  \noindent
  \begin{tabularx}{\linewidth}{@{}>{\raggedright\arraybackslash\bfseries}p{2.7cm}>{\raggedright\arraybackslash}X@{}}
    \toprule
    Probabilità & #3 \\
    Impatto     & #4 \\
    Livello     & #5 \\
    \midrule
    Descrizione & #6 \\
    \addlinespace
    Rilevamento & #7 \\
    \addlinespace
    Mitigazione & #8 \\
    \addlinespace
    Contingenza & #9 \\
    \bottomrule
  \end{tabularx}
  \medskip
}

\begin{document}

% ---------------------------------------------------------------------
% Contenuto del Documento
% ---------------------------------------------------------------------
\section{Introduzione}

\subsection{Scopo del documento}
Il presente documento ha lo scopo di pianificare, gestire e tracciare l'avanzamento delle attività necessarie allo sviluppo del progetto software. Il \textit{Piano di Progetto} fornisce una visione d'insieme sia a livello operativo che amministrativo, definendo la ripartizione dei compiti tra i membri del gruppo, le scadenze e l'impiego delle risorse.

Per pianificare il lavoro in modo chiaro e puntuale, garantendo il controllo dei costi e la gestione dei rischi, il documento tratta i seguenti argomenti:
\begin{itemize}
\item Identificazione dei rischi potenziali e pianificazione delle misure di mitigazione atte a minimizzarne l'impatto;
\item Definizione del modello di sviluppo adottato per la realizzazione del software;
\item Stime temporali ed economiche, attraverso un calendario di massima e un'analisi dei costi preventivati;
\item Pianificazione dettagliata delle attività da svolgere all'interno dei vari periodi di lavoro (sprint);
\item Rendicontazione economica e stesura dei consuntivi di periodo per monitorare il budget speso.
\end{itemize}

\subsection{Scopo del prodotto}
Il progetto, basato sul capitolato C2 denominato \textit{NEXUM Concierge}, ha come obiettivo l'introduzione di un assistente conversazionale all'interno della piattaforma aziendale NEXUM. Tale piattaforma, realizzata come \textit{Progressive Web App} (PWA), è pensata per avvicinare aziende e personale, semplificando le attività quotidiane legate alla comunicazione e alla gestione dei processi HR.

L'assistente permetterà agli utenti di cercare informazioni, comprendere le procedure aziendali e accedere ai servizi utilizzando direttamente il linguaggio naturale. Il \textit{Concierge} si occuperà di collegare la conoscenza aziendale alle funzionalità operative, mantenendo il contesto del dialogo e coordinando le interrogazioni alle fonti documentali o l'inoltro di operazioni tramite i servizi applicativi.
Per garantire un approccio graduale e verificabile, il dominio del \textit{Minimum Viable Product} (MVP) sarà focalizzato unicamente sull'area ferie e permessi.

\subsection{Glossario}
Al fine di evitare ambiguità e incomprensioni, e per massimizzare la chiarezza formale del testo, i termini tecnici, gli acronimi e le parole che necessitano di una definizione rigorosa sono raccolti all'interno del documento \textit{Glossario}. Tali termini sono contrassegnati all'interno dei testi di progetto tramite l'apposito pedice \ped{G}.

\subsection{Riferimenti}
% \subsubsection{Riferimenti normativi}
% \begin{itemize}
% \item \textbf{Capitolato d'appalto C2}: \textit{[NEXUM] BRD-AI-CONCIERGE-2026 - Business Requirements Document};
% \item \textbf{Norme di Progetto}: documento interno per la definizione delle regole e delle linee guida di gruppo.
% \end{itemize}

% \subsubsection{Riferimenti informativi}
% \begin{itemize}
% \item \textbf{Analisi dei Requisiti}: documento esterno \textit{Analisi dei Requisiti};
% \item Slide del corso di \textit{Ingegneria del Software} - Processi di ciclo di vita del Software;
% \item Slide del corso di \textit{Ingegneria del Software} - Modelli di sviluppo del Software;
% \item Slide del corso di \textit{Ingegneria del Software} - Gestione di Progetto.
% \end{itemize}

% =====================================================================
\section{Rischi attesi e mitigazione}\label{sec:rischi}
% =====================================================================

\subsection{Metodo}
L'analisi dei rischi individua in anticipo le situazioni che possono
compromettere il rispetto dei tempi, dei costi o della qualità, e stabilisce
per ognuna come accorgersene e come reagire. Il gruppo la svolge in
quattro passi, ripetuti a ogni \textit{sprint}:
\begin{enumerate}
  \item \textbf{identificazione}: i rischi sono ricavati dal capitolato C2,
        dall'incontro con il \glo{Proponente} del 2026-10-26 e
        dall'esperienza dei gruppi degli anni precedenti;
  \item \textbf{analisi}: a ogni rischio sono assegnati una probabilità e un
        impatto, da cui deriva il livello;
  \item \textbf{pianificazione}: per ogni rischio si definiscono gli
        indicatori che ne segnalano l'insorgere, le misure preventive e il
        piano di contingenza da attuare se si verifica;
  \item \textbf{controllo}: nella retrospettiva di fine \textit{sprint} il
        gruppo verifica quali rischi si sono presentati e quanto sono state
        efficaci le misure; il consuntivo di ogni periodo li riporta e
        questa sezione viene aggiornata di conseguenza.
\end{enumerate}

Ogni rischio ha un codice \texttt{R\textit{XN}}, dove \texttt{\textit{X}}
indica la categoria e \texttt{\textit{N}} è un numero progressivo:
\begin{itemize}
  \item \textbf{RT}: rischi tecnologici;
  \item \textbf{RR}: rischi legati ai \glo[Requisito]{requisiti} e al prodotto;
  \item \textbf{RO}: rischi organizzativi;
  \item \textbf{RP}: rischi personali e interpersonali.
\end{itemize}

La probabilità e l'impatto si misurano su tre livelli
(Tabella~\ref{tab:scale-rischi}); la loro combinazione dà il livello del
rischio (Tabella~\ref{tab:matrice-rischi}), che stabilisce la priorità con
cui il rischio viene sorvegliato.

\begin{table}[H]
  \centering
  \small
  \caption{Scale di probabilità e impatto\label{tab:scale-rischi}}
  \begin{tabularx}{\linewidth}{@{}l>{\raggedright\arraybackslash}X>{\raggedright\arraybackslash}X@{}}
    \toprule
    \intestazione{Valore} & \intestazione{Probabilità} & \intestazione{Impatto} \\
    \midrule
    Basso & poco probabile nell'intero progetto &
      ritardo inferiore a un giorno, nessun effetto sui documenti o sul prodotto \\
    Medio & può presentarsi una o due volte nel progetto &
      ritardo di alcuni giorni, recuperabile nello \textit{sprint} in corso o nel successivo \\
    Alto & atteso in più \textit{sprint} &
      mette a rischio una revisione (RTB o PB), il preventivo o un criterio di accettazione \\
    \bottomrule
  \end{tabularx}
\end{table}

\begin{table}[H]
  \centering
  \small
  \caption{Livello del rischio in base a probabilità e impatto\label{tab:matrice-rischi}}
  \begin{tabular}{@{}lccc@{}}
    \toprule
    \intestazione{Probabilità \textbackslash{} Impatto} & \intestazione{Basso} &
      \intestazione{Medio} & \intestazione{Alto} \\
    \midrule
    Alta  & Medio & Alto  & Alto  \\
    Media & Basso & Medio & Alto  \\
    Bassa & Basso & Basso & Medio \\
    \bottomrule
  \end{tabular}
\end{table}

Come strumenti di rilevamento il gruppo usa, oltre ai verbali e alle
retrospettive, il cruscotto interno alimentato dal GitHub Project: ore
svolte rispetto a quelle pianificate, consumo del preventivo per persona e
ruolo, obiettivi raggiunti e scostamento delle stime. 

% ---------------------------------------------------------------------
\subsection{Rischi tecnologici}

\rischio{RT1}{Inesperienza con le tecnologie del capitolato}{Alta}{Alto}{Alto}
{Il capitolato richiede tecnologie che il gruppo non ha mai usato in un
 progetto: modelli linguistici (\glo{LLM}) e la loro integrazione,
 \glo{RAG} con \glo{Embedding} e ricerca vettoriale, il protocollo
 \glo{MCP}, e lo \glo{Stack tecnologico} suggerito da Eggon (Ruby on Rails,
 Angular, RubyLLM, servizi AWS). A queste si aggiungono gli strumenti del
 modo di lavorare (\LaTeX{}, GitHub Actions). Lo studio necessario rallenta
 le prime attività e può portare a scelte poco consapevoli.}
{Stime superate in modo sistematico sulle attività tecniche (scostamento
 delle stime); attività ferme in \textit{In Progress} oltre la metà dello
 \textit{sprint}; domande ricorrenti sullo stesso argomento nelle riunioni.}
{Studio mirato a coppie nei primi \textit{sprint}, usando i \glo{PoC} come
 strumento di apprendimento.
 Brevi sessioni interne in cui chi ha studiato una tecnologia la presenta
 agli altri. Uso del materiale dei gruppi che hanno lavorato al capitolato
 negli anni precedenti, messo a disposizione da Eggon, e richiesta di
 incontri di approfondimento con il \glo{Proponente}, che si è reso
 disponibile.}
{Affiancamento del membro in difficoltà da parte di chi conosce già la
 tecnologia; ridistribuzione delle attività nello \textit{sprint}; incontro
 tecnico straordinario con Eggon; se necessario, adozione della soluzione
 più semplice tra quelle confrontate, concordata con il proponente.}

\rischio{RT2}{Qualità e affidabilità delle risposte del modello}{Alta}{Alto}{Alto}
{Un modello linguistico può dare risposte non supportate dalle fonti,
 citare la fonte sbagliata o non astenersi quando le informazioni non
 bastano; il suo comportamento non è deterministico e cambia con il
 \textit{prompt} e con il modello usato. Il capitolato chiede di dimostrare
 proprio citazioni, astensione e chiarimenti: una qualità insufficiente fa
 fallire i criteri di accettazione dell'\glo{MVP}.}
{Insieme di domande di prova con le risposte attese, costruito con Eggon a
 partire dal regolamento ferie e permessi e rieseguito a ogni modifica di
 \textit{prompt}, modello o catena di recupero delle fonti; metriche di
 qualità delle risposte (correttezza, citazioni corrette, astensioni
 corrette) definite nel Piano di Qualifica.}
{Risposte vincolate alle fonti recuperate, con citazione obbligatoria e
 astensione esplicita sotto una soglia di pertinenza; \textit{prompt}
 versionati nel repository come il codice; separazione netta tra
 conoscenza documentale e dati personali, che arrivano solo dai servizi
 di NEXUM: il modello non calcola né inventa dati. Confronto
 sperimentale tra documento intero nel contesto e \glo{RAG}, come chiesto
 dal capitolato, prima di scegliere l'approccio.}
{Restringere le domande supportate a quelle che superano le prove,
 rispondendo alle altre con l'astensione e il rimando alla fonte.}

\rischio{RT3}{Integrazione con i servizi di NEXUM}{Media}{Alto}{Alto}
{I dati personali (per esempio le ferie già godute) e le operazioni (la
 registrazione di una richiesta) passano dai servizi di NEXUM, secondo un
 contratto di integrazione definito da Eggon. Un contratto definito tardi,
 incompleto o modificato blocca lo sviluppo delle capacità operative, che
 sono il nucleo dell'MVP.}
{Stato del contratto di integrazione verificato a ogni incontro e
 riportato nei verbali esterni; attività bloccate in attesa di risposte
 del proponente.}
{Concordare il contratto entro la RTB; sviluppare su simulatori conformi
 al contratto, ammessi dal capitolato; isolare l'accesso ai servizi in un
 componente dedicato, così che il passaggio dal simulatore al servizio
 reale non tocchi il resto del sistema; prove automatiche sul rispetto del
 contratto.}
{Dimostrare il percorso completo dell'MVP sui simulatori conformi,
 documentando lo scostamento; ripianificare le attività dipendenti nello
 \textit{sprint} successivo.}

\rischio{RT4}{Scelta tecnologica inadeguata}{Media}{Medio}{Medio}
{Il capitolato non impone tecnologie e chiede scelte motivate da
 esperimenti. Una scelta che si rivela inadatta (prestazioni, curva di
 apprendimento, compatibilità) dopo l'avvio dello sviluppo costringe a
 rifare parte del lavoro.}
{Esiti dei PoC valutati rispetto a criteri di successo scritti prima di
 iniziare; problemi ricorrenti con una libreria o un servizio durante lo
 sviluppo.}
{Decisioni prese solo dopo esperimenti comparativi documentati; architettura a componenti con interfacce chiare (fornitore del
 modello, registro delle abilità MCP, accesso ai servizi), così che un
 componente si possa sostituire da solo.}
{Sostituzione del solo componente interessato dietro la sua interfaccia e
 ripianificazione dello \textit{sprint}.}

\rischio{RT5}{Malfunzionamenti di strumenti e infrastruttura}{Media}{Basso}{Basso}
{Servizi esterni del modo di lavorare (GitHub, GitHub Actions) in ritardo o
 non disponibili, errori nella compilazione automatica dei documenti,
 guasti dei computer personali o perdita di lavoro non salvato.}
{Esecuzioni automatiche in errore o ferme in coda; segnalazioni dei membri
 nei canali del gruppo.}
{Tutto il lavoro è versionato su GitHub con \textit{push} frequenti;
 documenti compilabili anche in locale; installazione degli strumenti
 della compilazione automatica conservata in cache.}
{Rilancio delle esecuzioni automatiche; compilazione locale; uso
 temporaneo di un altro computer o di un editor online.}

% ---------------------------------------------------------------------
\subsection{Rischi legati ai requisiti e al prodotto}

\rischio{RR1}{Ampliamento del perimetro}{Media}{Alto}{Alto}
{Il capitolato indica come possibilità ulteriori altri domini oltre a
 ferie e permessi, altri formati oltre al PDF, grafi di conoscenza e
 tecniche avanzate. Inseguirle sottrae ore al nucleo obbligatorio, cioè
 il percorso completo dalla domanda all'invio della richiesta.}
{Nuove richieste emerse negli incontri esterni; crescita dei requisiti
 desiderabili e opzionali nell'Analisi dei Requisiti; ore registrate su
 attività fuori dal perimetro dell'MVP.}
{Perimetro dell'MVP fissato e confermato con Eggon il 2026-10-26: solo
 ferie e permessi, solo lingua italiana, regolamento aziendale in PDF come
 prima fonte. Requisiti classificati in obbligatori, desiderabili e
 opzionali; ogni nuova richiesta viene stimata in ore prima di essere
 accettata e la decisione è registrata a verbale.}
{Rinvio delle richieste non essenziali a dopo il completamento dell'MVP;
 negoziazione con il proponente delle priorità.}

\rischio{RR2}{Requisiti incompresi o incompleti}{Media}{Alto}{Alto}
{I sedici requisiti di business del capitolato vanno tradotti in
 requisiti di sistema verificabili. Interpretazioni errate (per esempio
 sulla differenza tra richiesta registrata e approvata, sui due profili
 utente o sulla separazione tra aziende) emergono tardi e costano
 rifacimenti.}
{Domande aperte che restano senza risposta per più di un incontro;
 casi d'uso respinti o corretti in sede di validazione; requisiti senza
 una fonte nel capitolato o nei verbali.}
{Incontro settimanale fisso con Eggon, con domande preparate in anticipo;
 validazione progressiva dei casi d'uso con il proponente; tracciamento
 di ogni requisito verso la sua fonte (capitolato o verbale) nell'Analisi
 dei Requisiti.}
{Incontro straordinario di chiarimento; revisione dei requisiti
 interessati e ripianificazione delle attività che ne dipendono.}

\rischio{RR3}{Operazioni errate o duplicate}{Bassa}{Alto}{Medio}
{L'assistente può registrare una richiesta di ferie sbagliata (date
 ambigue, anno mancante), inviarla due volte dopo un errore di rete o
 presentare come approvata una richiesta solo registrata. Il capitolato
 lo considera un criterio di accettazione.}
{Prove di accettazione sul percorso completo con casi limite (date
 ambigue, annullamenti, errori dei servizi); registro delle operazioni
 eseguite.}
{Riepilogo completo e conferma esplicita prima di ogni invio; richiesta
 di chiarimento quando un dato è ambiguo; controllo dell'esito
 dell'operazione prima di ripeterla; messaggi che distinguono sempre
 richiesta registrata e approvata.}
{Disattivazione dell'invio, mantenendo la sola preparazione della
 richiesta, finché il difetto non è corretto.}

\rischio{RR4}{Sicurezza, riservatezza e autorizzazioni}{Bassa}{Alto}{Medio}
{Il sistema tratta dati personali del dipendente, distingue due profili
 (personale e amministrazione) e separa i dati delle aziende. Un messaggio
 o un documento costruito ad arte può indurre il modello a operazioni non
 autorizzate; il codice prodotto con strumenti di generazione, ammessi da
 Eggon, può introdurre vulnerabilità.}
{Prove specifiche su autorizzazioni e separazione dei dati; revisione di
 sicurezza del codice nella verifica, con riferimento alle linee guida
 OWASP indicate da Eggon.}
{Solo dati fittizi nella sperimentazione; autorizzazioni applicate dai
 servizi e non dal modello; abilità esposte tramite MCP con ingressi e
 uscite validati e privilegi minimi; il codice generato con strumenti
 automatici viene compreso e verificato come qualsiasi altro codice.}
{Blocco della funzionalità interessata, correzione con priorità massima e
 comunicazione al proponente.}

% ---------------------------------------------------------------------
\subsection{Rischi organizzativi}

\rischio{RO1}{Sovrapposizione con lezioni, esami e festività}{Alta}{Medio}{Alto}
{La sessione d'esami invernale (gennaio e febbraio) coincide con la
 preparazione della RTB e con l'avvio della PB, e le festività natalizie
 cadono negli \textit{sprint} 4 e 5. La disponibilità del gruppo cala
 proprio in periodi decisivi.}
{Indisponibilità comunicate da ciascuno all'inizio dello \textit{sprint};
 ore svolte inferiori a quelle pianificate nel cruscotto.}
{Calendario con ore ridotte negli \textit{sprint} delle festività
 (Sezione~\ref{sec:calendario}); lavoro documentale anticipato a novembre e
 dicembre; candidatura alla RTB a fine gennaio con margine; incontro
 settimanale con Eggon mantenuto anche nei periodi d'esame, come chiesto
 dal proponente.}
{Riassegnazione delle attività tra i membri disponibili; rinvio delle
 attività non critiche; aggiornamento del calendario di massima.}

\rischio{RO2}{Sottostima delle attività e ritardi}{Alta}{Alto}{Alto}
{Per tutti è il primo progetto di queste dimensioni: le stime tendono a
 essere ottimistiche e i ritardi si accumulano verso le revisioni.}
{Scostamento tra ore stimate ed effettive; obiettivi dello \textit{sprint}
 non raggiunti; attività riportate allo \textit{sprint} successivo.}
{Attività piccole, che stanno dentro uno \textit{sprint} (indicativamente
 non più di 8 ore); stime riviste in ogni retrospettiva alla luce dei
 consuntivi; priorità esplicite per gli obiettivi di ogni \textit{sprint};
 margine prima di ogni revisione.}
{Ripianificazione dello \textit{sprint}, riduzione degli obiettivi meno
 importanti e, se serve, spostamento della data di revisione comunicato
 per tempo al \glo{Committente}.}

\rischio{RO3}{Rotazione dei ruoli e passaggio di consegne}{Media}{Medio}{Medio}
{Ogni due settimane ogni membro cambia ruolo. Chi subentra può non
 conoscere lo stato delle attività, e il lavoro a cavallo tra due
 \textit{sprint} rischia di restare senza responsabile.}
{Attività senza assegnatario a inizio \textit{sprint}; domande ricorrenti
 sullo stato del lavoro nei primi giorni dopo il cambio.}
{Piano di rotazione fissato in anticipo per l'intero progetto; passaggio
 di consegne nella riunione di fine \textit{sprint}, riportato nel
 verbale; ogni attività descritta nella sua issue con lo stato aggiornato;
 procedure scritte nelle Norme di Progetto.}
{Affiancamento da parte del membro uscente nei primi giorni del nuovo
 \textit{sprint}.}

\rischio{RO4}{Superamento del preventivo}{Media}{Alto}{Alto}
{Il preventivo fissa 92 ore per persona, ripartite per ruolo. Consumare
 troppo presto le ore di un ruolo (per esempio il Progettista) o
 eccedere nei ruoli più costosi (Responsabile, 30\,\texteuro/h) lascia
 senza risorse le fasi finali.}
{Ore consumate e residue per persona e per ruolo nel cruscotto, aggiornato
 ogni notte; consumo del budget rispetto al tempo trascorso; consuntivo di
 ogni periodo.}
{Pianificazione degli \textit{sprint} coerente con il preventivo
 (Sezione~\ref{sec:impegno}); registrazione puntuale delle ore su ogni
 attività; controllo del consumo a ogni retrospettiva.}
{Riequilibrio dei ruoli negli \textit{sprint} successivi; revisione del
 preventivo a finire; riduzione del perimetro concordata con il
 proponente.}

\rischio{RO5}{Comunicazione discontinua con il proponente}{Bassa}{Alto}{Medio}
{Eggon ha indicato la continuità della comunicazione come rischio
 principale del progetto: incontri saltati o risposte tardive rallentano
 analisi e integrazione.}
{Incontri settimanali saltati; domande al proponente senza risposta da
 più di una settimana.}
{Incontro settimanale fisso di circa un'ora; canale rapido concordato per le comunicazioni brevi; verbale
 esterno di ogni incontro con decisioni e azioni; aggiornamento proattivo
 sulle scelte progettuali.}
{Incontro di recupero; comunicazione scritta delle decisioni in sospeso
 con una proposta del gruppo su cui il proponente possa rispondere.}

% ---------------------------------------------------------------------
\subsection{Rischi personali e interpersonali}

\rischio{RP1}{Indisponibilità di un membro}{Media}{Medio}{Medio}
{Malattie, impegni personali o, nel caso peggiore, l'abbandono del
 progetto riducono le ore disponibili e possono lasciare scoperte
 competenze che una sola persona conosce.}
{Comunicazioni di indisponibilità; ore registrate assenti per più giorni;
 attività assegnate ferme.}
{Nessuna competenza in mano a una sola persona: studio a coppie,
 documentazione delle scelte, attività descritte nelle issue; ore e
 attività riassegnabili grazie al tracciamento sul Project.}
{Ridistribuzione temporanea delle attività; in caso di abbandono,
 comunicazione immediata al committente e rinegoziazione del perimetro con
 il proponente.}

\rischio{RP2}{Divergenze nelle decisioni}{Media}{Basso}{Basso}
{Opinioni diverse su scelte tecniche o organizzative possono allungare le
 discussioni e rallentare le decisioni.}
{Lo stesso argomento discusso in più riunioni senza decisione.}
{Decisioni prese in riunione e registrate a verbale; per le scelte
 tecniche, confronto basato su esperimenti; se il consenso non si
 raggiunge, decide il Responsabile o si vota a maggioranza.}
{Mediazione del Responsabile; nei casi estremi, confronto con il
 proponente o con il committente.}

\rischio{RP3}{Comunicazione interna insufficiente}{Media}{Medio}{Medio}
{Informazioni disperse tra i canali del gruppo possono causare lavoro
 duplicato, decisioni non conosciute da tutti o attività dimenticate.}
{Due persone che lavorano alla stessa cosa; decisioni ignorate da parte
 del gruppo.}
{Canali e loro uso definiti nelle Norme di Progetto; ogni attività ha una
 issue, che è il riferimento per il suo stato; verbale dopo ogni
 riunione.}
{Riunione di riallineamento e aggiornamento delle issue interessate.}

% ---------------------------------------------------------------------
\subsection{Riepilogo}
La Tabella~\ref{tab:riepilogo-rischi} riassume i rischi in ordine di
livello. Quelli di livello alto sono riesaminati a ogni retrospettiva;
gli altri a ogni cambio di fase (RTB, PB).

\begin{table}[H]
  \centering
  \small
  \caption{Riepilogo dei rischi, dal livello più alto\label{tab:riepilogo-rischi}}
  \begin{tabularx}{\linewidth}{@{}l>{\raggedright\arraybackslash}Xccc@{}}
    \toprule
    \intestazione{Codice} & \intestazione{Rischio} & \intestazione{Probabilità} &
      \intestazione{Impatto} & \intestazione{Livello} \\
    \midrule
    RT1 & Inesperienza con le tecnologie del capitolato & Alta  & Alto  & Alto \\
    RT2 & Qualità e affidabilità delle risposte del modello & Alta & Alto & Alto \\
    RO2 & Sottostima delle attività e ritardi            & Alta  & Alto  & Alto \\
    RO1 & Sovrapposizione con lezioni, esami e festività  & Alta  & Medio & Alto \\
    RT3 & Integrazione con i servizi di NEXUM            & Media & Alto  & Alto \\
    RR1 & Ampliamento del perimetro                      & Media & Alto  & Alto \\
    RR2 & Requisiti incompresi o incompleti              & Media & Alto  & Alto \\
    RO4 & Superamento del preventivo                     & Media & Alto  & Alto \\
    \midrule
    RT4 & Scelta tecnologica inadeguata                  & Media & Medio & Medio \\
    RO3 & Rotazione dei ruoli e passaggio di consegne    & Media & Medio & Medio \\
    RP1 & Indisponibilità di un membro                   & Media & Medio & Medio \\
    RP3 & Comunicazione interna insufficiente            & Media & Medio & Medio \\
    RR3 & Operazioni errate o duplicate                  & Bassa & Alto  & Medio \\
    RR4 & Sicurezza, riservatezza e autorizzazioni       & Bassa & Alto  & Medio \\
    RO5 & Comunicazione discontinua con il proponente    & Bassa & Alto  & Medio \\
    \midrule
    RT5 & Malfunzionamenti di strumenti e infrastruttura & Media & Basso & Basso \\
    RP2 & Divergenze nelle decisioni                     & Media & Basso & Basso \\
    \bottomrule
  \end{tabularx}
\end{table}

% =====================================================================
% Qui andrà la sezione "Modello di sviluppo", come previsto
% dall'introduzione:
% \section{Modello di sviluppo}
% =====================================================================

% =====================================================================
\section{Calendario di massima}\label{sec:calendario}
% =====================================================================

Il progetto si articola in due revisioni di avanzamento: la
\textit{Requirements and Technology Baseline} (RTB), che fissa i requisiti
e le tecnologie dimostrandole con un \glo{PoC}, e la \textit{Product
Baseline} (PB), che presenta l'\glo{MVP} completo. La data della PB
coincide con l'impegno assunto nella Dichiarazione di Impegni
(Tabella~\ref{tab:revisioni}).

\begin{table}[H]
  \centering
  \small
  \caption{Date previste per le revisioni\label{tab:revisioni}}
  \begin{tabular}{@{}llc@{}}
    \toprule
    \intestazione{Revisione} & \intestazione{Contenuto} & \intestazione{Candidatura prevista} \\
    \midrule
    RTB & requisiti, tecnologie scelte e PoC & 2027-01-25 \\
    PB  & MVP completo e verificato          & 2027-03-23 \\
    \bottomrule
  \end{tabular}
\end{table}

Il lavoro è diviso in dieci \textit{sprint} di due settimane, da lunedì a
domenica, allineati alla rotazione dei ruoli
(Tabella~\ref{tab:calendario}). I primi sei portano alla RTB, gli ultimi
quattro alla PB. Gli \textit{sprint} 4 e 5 comprendono le festività
natalizie e prevedono circa un quarto di ore in meno degli altri; la RTB
cade all'inizio della sessione d'esami invernale, che il calendario tiene
in conto (rischio RO1, sezione~\ref{rischio:RO1}).

\begin{table}[H]
  \centering
  \small
  \caption{Calendario di massima degli \textit{sprint}\label{tab:calendario}}
  \begin{tabularx}{\linewidth}{@{}cllc>{\raggedright\arraybackslash}X@{}}
    \toprule
    \intestazione{Sprint} & \intestazione{Inizio} & \intestazione{Fine} &
      \intestazione{Fase} & \intestazione{Obiettivo principale} \\
    \midrule
    1  & 2026-11-02 & 2026-11-15 & RTB & avvio: pianificazione, norme, infrastruttura, primo incontro di analisi \\
    2  & 2026-11-16 & 2026-11-29 & RTB & casi d'uso delle quattro capacità; primi esperimenti con i modelli \\
    3  & 2026-11-30 & 2026-12-13 & RTB & requisiti completi e tracciati; esperimento sulla conoscenza documentale \\
    4  & 2026-12-14 & 2026-12-27 & RTB & PoC delle abilità MCP su servizi simulati (festività) \\
    5  & 2026-12-28 & 2027-01-10 & RTB & PoC del percorso completo; bozza della \textit{Technology Baseline} (festività) \\
    6  & 2027-01-11 & 2027-01-24 & RTB & verifica e approvazione dei documenti; preparazione della RTB \\
    \midrule
    7  & 2027-01-25 & 2027-02-07 & PB  & revisione RTB e correzioni; architettura dell'MVP; avvio dello sviluppo \\
    8  & 2027-02-08 & 2027-02-21 & PB  & MVP: conoscenza documentale, gestione delle fonti, conversazione \\
    9  & 2027-02-22 & 2027-03-07 & PB  & MVP: dati personali, richiesta con conferma, autorizzazioni \\
    10 & 2027-03-08 & 2027-03-21 & PB  & collaudo sui criteri di accettazione; documenti finali; preparazione della PB \\
    \bottomrule
  \end{tabularx}
\end{table}

Le date sono una stima iniziale: a ogni retrospettiva il consuntivo del
periodo le confronta con l'avanzamento reale e, se serve, questa sezione
viene aggiornata e le modifiche comunicate al committente.

% =====================================================================
\section{Impegno orario previsto e stima dei costi}\label{sec:impegno}
% =====================================================================

Il preventivo riprende quello della Dichiarazione di Impegni: 552 ore,
92 per ogni membro, per un costo complessivo di 11.250,00\,\texteuro.
Ogni membro ricopre tutti i ruoli con la stessa ripartizione oraria, così
che il carico e le opportunità di apprendimento siano uguali per tutti
(Tabella~\ref{tab:ore-membro}). La Tabella~\ref{tab:costi-ruolo} riporta
ore e costi per ruolo.

\begin{table}[H]
  \centering
  \small
  \caption{Ore previste per membro e per ruolo\label{tab:ore-membro}}
  \begin{tabular}{@{}lrrrrrrr@{}}
    \toprule
    \intestazione{Membro} & \intestazione{Re} & \intestazione{Am} & \intestazione{An} &
      \intestazione{Pt} & \intestazione{Pr} & \intestazione{Ve} & \intestazione{Totale} \\
    \midrule
    \nome{biasini}    & 9 & 12 & 13 & 17 & 22 & 19 & 92 \\
    \nome{cescobolla} & 9 & 12 & 13 & 17 & 22 & 19 & 92 \\
    \nome{corro}      & 9 & 12 & 13 & 17 & 22 & 19 & 92 \\
    \nome{guglielmo}  & 9 & 12 & 13 & 17 & 22 & 19 & 92 \\
    \nome{mazzurana}  & 9 & 12 & 13 & 17 & 22 & 19 & 92 \\
    \nome{spagnoli}   & 9 & 12 & 13 & 17 & 22 & 19 & 92 \\
    \midrule
    \textbf{Totale}   & 54 & 72 & 78 & 102 & 132 & 114 & 552 \\
    \bottomrule
  \end{tabular}

  \smallskip
  {\footnotesize Re = Responsabile, Am = Amministratore, An = Analista,
   Pt = Progettista, Pr = Programmatore, Ve = Verificatore.}
\end{table}

\begin{table}[H]
  \centering
  \small
  \caption{Ore e costi previsti per ruolo\label{tab:costi-ruolo}}
  \begin{tabular}{@{}lrrrr@{}}
    \toprule
    \intestazione{Ruolo} & \intestazione{Ore} & \intestazione{Quota ore} &
      \intestazione{Costo orario} & \intestazione{Costo} \\
    \midrule
    Responsabile   &  54 &  9,8\,\% & 30,00\,\texteuro & 1.620,00\,\texteuro \\
    Amministratore &  72 & 13,0\,\% & 20,00\,\texteuro & 1.440,00\,\texteuro \\
    Analista       &  78 & 14,1\,\% & 25,00\,\texteuro & 1.950,00\,\texteuro \\
    Progettista    & 102 & 18,5\,\% & 25,00\,\texteuro & 2.550,00\,\texteuro \\
    Programmatore  & 132 & 23,9\,\% & 15,00\,\texteuro & 1.980,00\,\texteuro \\
    Verificatore   & 114 & 20,7\,\% & 15,00\,\texteuro & 1.710,00\,\texteuro \\
    \midrule
    \textbf{Totale} & \textbf{552} & 100\,\% & & \textbf{11.250,00\,\texteuro} \\
    \bottomrule
  \end{tabular}
\end{table}

La ripartizione tra le due fasi segue la natura del lavoro
(Tabella~\ref{tab:ore-fase}). Verso la RTB prevalgono analisi,
amministrazione e verifica dei documenti, oltre alla prima progettazione e
al codice del PoC; verso la PB prevalgono progettazione di dettaglio,
programmazione e verifica del prodotto. Poco più della metà delle ore è
destinata alla RTB, perché in questa fase si producono quasi tutti i
documenti e si sceglie lo stack tecnologico.

\begin{table}[H]
  \centering
  \small
  \caption{Ore e costi previsti per fase e ruolo\label{tab:ore-fase}}
  \begin{tabular}{@{}lrrrrrrrr@{}}
    \toprule
    \intestazione{Fase} & \intestazione{Re} & \intestazione{Am} & \intestazione{An} &
      \intestazione{Pt} & \intestazione{Pr} & \intestazione{Ve} &
      \intestazione{Ore} & \intestazione{Costo} \\
    \midrule
    RTB (\textit{sprint} 1--6)  & 31,5 & 48 & 65 & 47 & 56 & 64 & 311,5 & 6.505,00\,\texteuro \\
    PB (\textit{sprint} 7--10)  & 22,5 & 24 & 13 & 55 & 76 & 50 & 240,5 & 4.745,00\,\texteuro \\
    \midrule
    \textbf{Totale} & 54 & 72 & 78 & 102 & 132 & 114 & 552 & 11.250,00\,\texteuro \\
    \bottomrule
  \end{tabular}
\end{table}

La suddivisione di dettaglio per \textit{sprint}, membro e ruolo è
riportata nella pianificazione di ogni periodo. Il consumo effettivo è
tracciato con le ore registrate su ogni attività e confrontato con il
preventivo nel consuntivo di ogni periodo (rischio RO4, sezione~\ref{rischio:RO4}).

% =====================================================================
\section{Pianificazione}\label{sec:pianificazione}
% =====================================================================

\subsection{Introduzione}
Questa sezione raccoglie, per ogni \textit{sprint}, la pianificazione
decisa nella riunione di apertura e il resoconto scritto alla chiusura.
Ogni \textit{sprint} è descritto con la stessa struttura:
\begin{itemize}
  \item \textbf{Panoramica e obiettivi}: il contesto del periodo e i
        risultati attesi, ognuno tracciato da una o più \textit{issue};
  \item \textbf{Rischi attesi}: i rischi della sezione~\ref{sec:rischi}
        più probabili nel periodo e come vengono mitigati;
  \item \textbf{Ruoli assegnati}: il ruolo di ogni membro, uno per
        persona, secondo la rotazione;
  \item \textbf{Preventivo}: ore e costi pianificati per membro e per ruolo;
  \item \textbf{Consuntivo}: ore e costi effettivi, con lo scarto rispetto
        al preventivo;
  \item \textbf{Risorse rimanenti}: ore e budget ancora disponibili dopo
        lo \textit{sprint};
  \item \textbf{Retrospettiva}: rischi incontrati, cosa ha funzionato,
        cosa migliorare e come cambia la pianificazione successiva.
\end{itemize}

Le ore del consuntivo non sono stimate a posteriori: ogni membro le
registra durante il lavoro sulla \textit{issue} dell'attività svolta, con
il comando \texttt{/ore} descritto nelle \textit{Norme di Progetto}. Le
registrazioni vengono sommate automaticamente per persona, ruolo e
\textit{sprint}, e da queste derivano le tabelle di consuntivo. Ruoli,
ore pianificate e obiettivi di ogni \textit{sprint} sono gli stessi
riportati nel file di pianificazione del gruppo, da cui il cruscotto
confronta ogni giorno ore pianificate ed effettive.

Nelle tabelle i ruoli sono abbreviati come segue: Re = Responsabile,
Am = Amministratore, An = Analista, Pt = Progettista, Pr = Programmatore,
Ve = Verificatore. Le tariffe orarie sono quelle della
Tabella~\ref{tab:costi-ruolo}.

% =====================================================================
\subsection{Requirements and Technology Baseline}\label{sec:pianificazione-rtb}
% =====================================================================

% ---------------------------------------------------------------------
\subsubsection{Sprint 1: dal 2026-11-02 al 2026-11-15}\label{sec:sprint-1}
% ---------------------------------------------------------------------

\subsubsubsection{Panoramica e obiettivi}
Il primo \textit{sprint} si apre con la presentazione del Diario di Bordo
del 2026-11-02 e con l'aggiudicazione del capitolato C2
\textit{NEXUM Concierge}, presa in esame dal gruppo nella riunione interna
del 2026-11-03. Il periodo ha un carattere soprattutto organizzativo: il
gruppo imposta i processi e gli strumenti che userà per tutto il progetto,
avvia i documenti della RTB e apre il dialogo con la proponente Eggon, da
cui dipendono l'Analisi dei Requisiti e le scelte tecnologiche dei periodi
successivi. Gli obiettivi fissati sono:
\begin{itemize}
  \item \textbf{Piano di Progetto}: prima stesura di rischi attesi e
        mitigazione, calendario di massima, impegno orario e
        pianificazione del primo \textit{sprint};
  \item \textbf{Norme di Progetto}: prima stesura, con i processi di
        gestione, di documentazione e di verifica già adottati dal gruppo;
  \item \textbf{Infrastruttura di gestione}: organizzazione del GitHub
        Project con i campi per \textit{sprint}, ruolo e ore;
        collegamento automatico tra \textit{issue}, \textit{pull request}
        e Project; registrazione delle ore con il comando \texttt{/ore};
        cruscotto di avanzamento; creazione del repository per i
        \glo{PoC};
  \item \textbf{Glossario}: riorganizzazione per lettere e inserimento dei
        primi termini del dominio del capitolato;
  \item \textbf{Rapporti con la proponente}: primo incontro con Eggon per
        concordare canali e cadenza di comunicazione e per chiarire il
        perimetro dell'MVP;
  \item \textbf{Analisi dei Requisiti}: studio del capitolato,
        individuazione degli attori e delle quattro capacità del
        \textit{Concierge}, prima bozza dei casi d'uso;
  \item \textbf{Studio delle tecnologie}: approfondimento delle tecnologie
        indicate nel capitolato e preparazione dell'ambiente per i primi
        esperimenti;
  \item \textbf{Verifica}: revisione di ogni documento e verbale prodotto
        nel periodo prima della sua approvazione.
\end{itemize}

\subsubsubsection{Rischi attesi}
Per il primo \textit{sprint} il gruppo prevede i seguenti rischi:
\begin{itemize}
  \item \textbf{RT1 -- Inesperienza con le tecnologie del capitolato}
        (sezione~\ref{rischio:RT1}): nessun membro ha ancora lavorato con
        le tecnologie richieste. Lo studio è distribuito tra i componenti del gruppo e i risultati vengono condivisi nelle riunioni interne;
  \item \textbf{RR2 -- Requisiti incompresi o incompleti}
        (sezione~\ref{rischio:RR2}): il capitolato lascia aperti diversi
        aspetti dell'integrazione con NEXUM. I dubbi vengono raccolti in
        un elenco di domande da portare al primo incontro con Eggon;
  \item \textbf{RO5 -- Comunicazione discontinua con il proponente}
        (sezione~\ref{rischio:RO5}): il canale con Eggon è ancora da
        stabilire. Il Responsabile concorda nel primo incontro modalità e
        frequenza dei contatti;
\end{itemize}

\subsubsubsection{Ruoli assegnati}
I ruoli del primo \textit{sprint} sono riportati nella
Tabella~\ref{tab:ruoli-sprint-1}. I ruoli seguono la rotazione pianificata, che
garantisce a ogni membro le ore previste per ciascun ruolo entro la fine
del progetto.

\begin{table}[H]
  \centering
  \small
  \caption{Sprint 1 -- Ruoli assegnati\label{tab:ruoli-sprint-1}}
  \begin{tabular}{@{}llr@{}}
    \toprule
    \intestazione{Ruolo} & \intestazione{Membro} & \intestazione{Ore pianificate} \\
    \midrule
    Responsabile   & \nome{cescobolla} & 4,5 \\
    Amministratore & \nome{biasini}    & 12 \\
    Analista       & \nome{guglielmo}  & 13 \\
                   & \nome{mazzurana}  & 6 \\
    Progettista    & \nome{spagnoli}   & 8 \\
    Programmatore  & --                & -- \\
    Verificatore   & \nome{corro}      & 9,5 \\
    \bottomrule
  \end{tabular}
\end{table}

Non è previsto un Programmatore: nel primo periodo non si scrive ancora
codice di prodotto, e le prime prove tecniche rientrano nello studio
svolto dal Progettista. Il peso maggiore è sull'analisi (19 ore), perché
l'Analisi dei Requisiti è il documento su cui poggia tutta la RTB.

\subsubsubsection{Preventivo}
La Tabella~\ref{tab:preventivo-sprint-1} riporta le ore pianificate per
membro e ruolo; la Tabella~\ref{tab:preventivo-economico-sprint-1} le
riassume per ruolo con il relativo costo.

\begin{table}[H]
  \centering
  \small
  \caption{Sprint 1 -- Preventivo orario per membro e ruolo\label{tab:preventivo-sprint-1}}
  \begin{tabular}{@{}lrrrrrrr@{}}
    \toprule
    \intestazione{Membro} & \intestazione{Re} & \intestazione{Am} & \intestazione{An} &
      \intestazione{Pt} & \intestazione{Pr} & \intestazione{Ve} & \intestazione{Totale} \\
    \midrule
    \nome{biasini}    & --  & 12 & -- & -- & -- & --  & 12 \\
    \nome{cescobolla} & 4,5 & -- & -- & -- & -- & --  & 4,5 \\
    \nome{corro}      & --  & -- & -- & -- & -- & 9,5 & 9,5 \\
    \nome{guglielmo}  & --  & -- & 13 & -- & -- & --  & 13 \\
    \nome{mazzurana}  & --  & -- & 6  & -- & -- & --  & 6 \\
    \nome{spagnoli}   & --  & -- & -- & 8  & -- & --  & 8 \\
    \midrule
    \textbf{Totale ore}         & 4,5    & 12     & 19     & 8      & -- & 9,5    & 53 \\
    \textbf{Costo (\texteuro)}  & 135,00 & 240,00 & 475,00 & 200,00 & -- & 142,50 & 1.192,50 \\
    \bottomrule
  \end{tabular}
\end{table}

\begin{table}[H]
  \centering
  \small
  \caption{Sprint 1 -- Preventivo economico per ruolo\label{tab:preventivo-economico-sprint-1}}
  \begin{tabular}{@{}lrrr@{}}
    \toprule
    \intestazione{Ruolo} & \intestazione{Ore} & \intestazione{Quota ore} & \intestazione{Costo} \\
    \midrule
    Responsabile   & 4,5 &  8,5\,\% & 135,00\,\texteuro \\
    Amministratore & 12  & 22,6\,\% & 240,00\,\texteuro \\
    Analista       & 19  & 35,8\,\% & 475,00\,\texteuro \\
    Progettista    & 8   & 15,1\,\% & 200,00\,\texteuro \\
    Programmatore  & --  & --       & -- \\
    Verificatore   & 9,5 & 17,9\,\% & 142,50\,\texteuro \\
    \midrule
    \textbf{Totale} & \textbf{53} & 100\,\% & \textbf{1.192,50\,\texteuro} \\
    \bottomrule
  \end{tabular}
\end{table}

\begin{figure}[H]
  \centering
  \begin{tikzpicture}[x=0.32cm, y=0.62cm, font=\small]
    % una barra per ruolo, lunga quanto le ore preventivate
    \foreach \ruolo/\ore/\testo [count=\i] in
      {Responsabile/4.5/4{,}5, Amministratore/12/12, Analista/19/19,
       Progettista/8/8, Programmatore/0/0, Verificatore/9.5/9{,}5} {
      \node[anchor=east] at (-0.4, -\i) {\ruolo};
      \ifdim \ore pt > 0pt
        \fill[gptaccento] (0, -\i-0.32) rectangle (\ore, -\i+0.32);
      \fi
      \node[anchor=west, text=gptinchiostro] at (\ore+0.3, -\i) {\testo\ h};
    }
    \draw[gptgrigio] (0, -0.4) -- (0, -6.6);
  \end{tikzpicture}
  \caption{Sprint 1 -- Ore preventivate per ruolo\label{fig:preventivo-sprint-1}}
\end{figure}

\subsubsubsection{Consuntivo}
\textit{Da compilare alla chiusura dello sprint (2026-11-15).}

% ---------------------------------------------------------------------
% COME COMPILARE IL CONSUNTIVO
% Le ore effettive si leggono dal cruscotto (repository .github-private):
% il file tabelle/consuntivo-sprint-1.tex contiene già questa tabella,
% con lo scarto dal preventivo tra parentesi. Si può incollare al posto
% dello scheletro qui sotto oppure includerlo con
%   \input{tabelle/consuntivo-sprint-1}
% Togliete la riga "Da compilare..." e il simbolo % davanti alle righe.
% Nelle celle: ore effettive e, se diverse, lo scarto, per esempio
%   13 \scriptsize(+1)   oppure   5 \scriptsize($-$1)
% ---------------------------------------------------------------------
% La Tabella~\ref{tab:consuntivo-sprint-1} riporta le ore effettivamente
% registrate per membro e ruolo, con lo scarto dal preventivo tra
% parentesi; la Tabella~\ref{tab:consuntivo-economico-sprint-1} confronta
% ore e costi per ruolo con il preventivo.
%
% \begin{table}[H]
%   \centering
%   \small
%   \caption{Sprint 1 -- Consuntivo orario per membro e ruolo\label{tab:consuntivo-sprint-1}}
%   \begin{tabular}{@{}lrrrrrrr@{}}
%     \toprule
%     \intestazione{Membro} & \intestazione{Re} & \intestazione{Am} & \intestazione{An} &
%       \intestazione{Pt} & \intestazione{Pr} & \intestazione{Ve} & \intestazione{Totale} \\
%     \midrule
%     \nome{biasini}    & --  & ... & --  & --  & -- & --  & ... \\
%     \nome{cescobolla} & ... & --  & --  & --  & -- & --  & ... \\
%     \nome{corro}      & --  & --  & --  & --  & -- & ... & ... \\
%     \nome{guglielmo}  & --  & --  & ... & --  & -- & --  & ... \\
%     \nome{mazzurana}  & --  & --  & ... & --  & -- & --  & ... \\
%     \nome{spagnoli}   & --  & --  & --  & ... & -- & --  & ... \\
%     \midrule
%     \textbf{Totale ore}        & ... & ... & ... & ... & -- & ... & ... \\
%     \textbf{Costo (\texteuro)} & ... & ... & ... & ... & -- & ... & ... \\
%     \bottomrule
%   \end{tabular}
% \end{table}
%
% \begin{table}[H]
%   \centering
%   \small
%   \caption{Sprint 1 -- Consuntivo economico per ruolo\label{tab:consuntivo-economico-sprint-1}}
%   \begin{tabular}{@{}lrrrr@{}}
%     \toprule
%     \intestazione{Ruolo} & \intestazione{Ore} & \intestazione{Scarto ore} &
%       \intestazione{Costo} & \intestazione{Scarto costo} \\
%     \midrule
%     Responsabile   & ... & ... & ...\,\texteuro & ...\,\texteuro \\
%     Amministratore & ... & ... & ...\,\texteuro & ...\,\texteuro \\
%     Analista       & ... & ... & ...\,\texteuro & ...\,\texteuro \\
%     Progettista    & ... & ... & ...\,\texteuro & ...\,\texteuro \\
%     Programmatore  & --  & --  & --             & --             \\
%     Verificatore   & ... & ... & ...\,\texteuro & ...\,\texteuro \\
%     \midrule
%     \textbf{Totale} & ... & ... & ...\,\texteuro & ...\,\texteuro \\
%     \bottomrule
%   \end{tabular}
%
%   \smallskip
%   {\footnotesize Scarto = consuntivo $-$ preventivo: un valore positivo
%    indica più ore o più spesa del previsto.}
% \end{table}

\subsubsubsection{Risorse rimanenti}
\textit{Da compilare alla chiusura dello sprint (2026-11-15).}

% ---------------------------------------------------------------------
% Ore e budget rimanenti = preventivo totale (Tabella tab:costi-ruolo)
% meno le ore consumate dall'inizio del progetto. Il cruscotto le
% calcola nel file tabelle/riepilogo-ruoli.tex (colonne "res.").
% Partenza: 54 / 72 / 78 / 102 / 132 / 114 ore = 552 ore, 11.250,00 euro.
% ---------------------------------------------------------------------
% \begin{table}[H]
%   \centering
%   \small
%   \caption{Sprint 1 -- Ore e budget rimanenti per ruolo\label{tab:rimanenti-sprint-1}}
%   \begin{tabular}{@{}lrrrr@{}}
%     \toprule
%     \intestazione{Ruolo} & \intestazione{Ore consumate} & \intestazione{Ore rimanenti} &
%       \intestazione{Costo consumato} & \intestazione{Budget rimanente} \\
%     \midrule
%     Responsabile   & ... & ... & ...\,\texteuro & ...\,\texteuro \\
%     Amministratore & ... & ... & ...\,\texteuro & ...\,\texteuro \\
%     Analista       & ... & ... & ...\,\texteuro & ...\,\texteuro \\
%     Progettista    & ... & ... & ...\,\texteuro & ...\,\texteuro \\
%     Programmatore  & 0   & 132 & 0,00\,\texteuro & 1.980,00\,\texteuro \\
%     Verificatore   & ... & ... & ...\,\texteuro & ...\,\texteuro \\
%     \midrule
%     \textbf{Totale} & ... & ... & ...\,\texteuro & ...\,\texteuro \\
%     \bottomrule
%   \end{tabular}
% \end{table}

\subsubsubsection{Retrospettiva}
\textit{Da compilare dopo la riunione di retrospettiva di fine sprint.}

% ---------------------------------------------------------------------
% Traccia (togliete i % e completate):
%
% \textbf{Rischi incontrati.} Quali dei rischi attesi si sono verificati,
% con che effetto, e se la mitigazione prevista ha funzionato. Eventuali
% rischi non previsti (in quel caso aggiungeteli alla sezione dei rischi).
%
% \textbf{Cosa ha funzionato.}
% \begin{itemize}
%   \item ...
% \end{itemize}
%
% \textbf{Cosa migliorare.}
% \begin{itemize}
%   \item ...
% \end{itemize}
%
% \textbf{Obiettivi.} Quali obiettivi sono stati raggiunti e quali
% passano allo sprint 2, con il motivo.
%
% \textbf{Aggiornamento della pianificazione.} Le azioni correttive
% decise per lo sprint 2 e le eventuali modifiche a ruoli o ore
% (in quel caso aggiornate anche il file di pianificazione).
% ---------------------------------------------------------------------

\end{document}